% ================================================================= 7. HỆ THỐNG AGENT
\section{Khảo sát ý tưởng hệ thống nhiều mô hình}
\label{sec:agent}

Một hướng nhóm cân nhắc là: thay vì một mô hình duy nhất, xây một hệ thống nhiều tác tử, mỗi nhiệm vụ
con dùng mô hình tối ưu riêng cho nó. Mục này là kết quả khảo sát nghiêm túc ý tưởng đó, dựa trên bằng
chứng từ 30 công trình đã đọc và từ số đo của chính nhóm.

\subsection{Nhận xét đầu tiên: pipeline hiện tại đã là hệ dị thể}

Pipeline của nhóm hiện gồm bốn khối dùng bốn loại công cụ khác nhau: lọc Butterworth (xử lý tín hiệu
cổ điển), mẫu trung vị kết hợp bình phương tối thiểu (thống kê cổ điển), mạng tích chập giãn nở (học sâu),
và lấy đỉnh có thời gian trơ (luật xác định). Đây \textbf{đã là} nguyên tắc ``mỗi việc một công cụ tối ưu''.
Thêm một tầng điều phối lên trên chỉ đổi tên gọi chứ không thêm năng lực.

\subsection{Phân rã bảy nhiệm vụ con và mô hình tối ưu cho từng nhiệm vụ}

\begin{table}[H]\centering\small
\caption{Bảy nhiệm vụ con của bài toán và mô hình tối ưu theo bằng chứng từ y văn cùng số đo của nhóm.}
\label{tab:nhiemvu}
\begin{tabular}{@{}C{0.8cm} L{3.4cm} L{5.0cm} L{4.6cm}@{}}
\toprule
\textbf{\#} & \textbf{Nhiệm vụ con} & \textbf{Mô hình tối ưu theo bằng chứng} & \textbf{Trạng thái} \\
\midrule
1 & Đánh giá chất lượng tín hiệu & Còn bỏ ngỏ. Heuristic phổ sụp ngoài miền (\S\ref{sec:chonkenh}) &
\xau{Đây chính là C3} \\
2 & Dò QRS mẹ & Luật cổ điển kiểu Pan--Tompkins & Đã có, đủ tốt \\
3 & Khử ECG mẹ & \textbf{Còn tranh chấp thật}: trừ mẫu, lọc Kalman mở rộng, hay mạng tái tạo &
Cơ hội cải tiến \\
4 & Dò QRS thai & Đã bão hoà: đổi kiến trúc chỉ $+0{,}41$ điểm, $p = 0{,}7012$ & Đã có, khó cải thiện \\
5 & Sửa lỗi nhịp, nội suy khoảng mất & Làm mượt khoảng RR bằng luật sinh lý & Chưa làm \\
6 & Ước lượng nhịp tim thai & Tính trực tiếp từ khoảng RR & Đã có \\
7 & Quyết định từ chối trả lời & Chưa có tiền lệ cho tín hiệu thai & \xau{Đây chính là C3} \\
\bottomrule
\end{tabular}
\end{table}

Trong bảy nhiệm vụ, \textbf{năm} có lời giải tốt nhất là giải pháp cổ điển xác định, \textbf{một} đã bão
hoà, và chỉ \textbf{một} (khử ECG mẹ) còn tranh chấp thật sự về mô hình.

\subsection{Y văn nói gì về hướng nhiều mô hình}

Trong 30 công trình đã đọc, \textbf{không có công trình nào} dùng tác tử ngôn ngữ lớn hoặc kiến trúc
hỗn hợp chuyên gia có cổng học được cho bài toán điện tim thai. Nhưng \textbf{có} tiền lệ rõ ràng cho
hợp nhất và nối tầng:

\begin{table}[H]\centering\small
\caption{Tiền lệ hợp nhất và nối tầng trong y văn, kèm mức lợi ích công bố.}
\label{tab:ensemble}
\begin{tabular}{@{}L{5.4cm} L{4.4cm} L{4.6cm}@{}}
\toprule
\textbf{Công trình} & \textbf{Cơ chế} & \textbf{Mức lợi ích công bố} \\
\midrule
Behar và cộng sự 2014 (FUSE) & Hợp nhất nhiều bộ trích xuất & khoảng $+3{,}0$ điểm F1 so với phương
pháp đơn tốt nhất \\
Shokouhmand \& Tavassolian 2023 & Nối tầng khối khử nhiễu trước tách nguồn & $+22{,}45$ điểm \\
Asadi và cộng sự 2025 & Nối tầng hai mạng U-Net & $+1{,}1$ đến $+1{,}6$ điểm độ nhạy \\
Behar và cộng sự 2014 (cổng bSQI) & \textbf{Cổng từ chối} theo chất lượng & $+0{,}95$ đến $+1{,}11$ điểm
khi chỉ loại 3,6\% tín hiệu \\
\bottomrule
\end{tabular}
\end{table}

\begin{ghichu}[title={Kết luận quan trọng nhất của mục này}]
Giá trị của ``tác tử'' trong y văn điện tim thai nằm ở \textbf{quyền chọn} và \textbf{quyền từ chối},
không nằm ở việc luân chuyển giữa nhiều mô hình dò. Điều này trùng khít với đóng góp C3 mà nhóm đã đề xuất.
\end{ghichu}

\subsection{Còn bao nhiêu dư địa để cải tiến}

Trong miền, nhóm đang ở 99,21 (kênh mù nhãn) và 97,45 (trung bình bốn kênh). Trần là 100. Nghĩa là dư địa
còn lại chỉ \textbf{0,79} đến \textbf{2,55} điểm. Trong khi đó, ngưỡng phân giải thống kê thực nghiệm của
nhóm nằm đâu đó giữa 0,41 điểm (không phân biệt được, $p = 0{,}70$) và 4,53 điểm (phân biệt rõ,
$p = 0{,}0000$). Suy ra: \textbf{bất kỳ thành phần nào hứa hẹn dưới khoảng 0,8 điểm đều không chứng minh
được với $n = 5$}, nên không đáng xây.

Chỗ sụp thật sự là ngoài miền: 77,34 với phân bố lưỡng cực. Đó mới là nơi đáng đầu tư.

\subsection{Rủi ro và khuyến nghị}

\begin{canhbao}[title={Rủi ro lớn nhất không phải tốc độ}]
Rủi ro lớn nhất của một hệ bảy tác tử là nó phá hỏng \textbf{hai tài sản duy nhất} khiến đề tài này có
cửa đăng Q1:
\begin{enumerate}[leftmargin=1.5em,itemsep=2pt]
\item \textbf{Tính đơn giản tái lập được}: 113.481 tham số, 0,48\,MB, 4,35\,ms. Rất hiếm trong 30 công
trình đã đọc, và là điểm mạnh thật khi phản biện.
\item \textbf{Tính chặt chẽ thống kê trên $n = 5$ sản phụ.} Mỗi thành phần học được thêm vào là một chỗ
mới cho rò rỉ dữ liệu ẩn nấp và một siêu tham số mới phải biện minh.
\end{enumerate}
Y văn có ví dụ cụ thể: FUSE với mười hai siêu tham số đã buộc chính tác giả phải báo cáo F1 trong mẫu.
\end{canhbao}

\textbf{Khuyến nghị của nhóm, theo thứ tự ưu tiên:}

\begin{enumerate}[leftmargin=1.8em,itemsep=4pt]
\item \textbf{Xây cổng từ chối dựa trên fSQI}, báo cáo bằng đường cong rủi ro--độ phủ kèm đường cơ sở
cổng ngẫu nhiên. Đây là C3.
\item \textbf{Làm mượt khoảng RR} bằng luật sinh lý, kèm bắt buộc so với đường cơ sở nhịp hằng số.
\item \textbf{Thí nghiệm chẩn đoán} vì sao ba bản ghi CinC sụp đổ.
\item Chỉ khi ba việc trên xong mới cân nhắc bộ chọn dải thông thích nghi --- dạng ``tác tử'' duy nhất
có căn cứ, vì nhóm đã chứng minh dải thông là biến quan trọng nhất.
\end{enumerate}

\textbf{Không} xây hệ bảy tác tử. \textbf{Không} dùng mô hình ngôn ngữ lớn điều phối. Cả hai đều thêm độ
phức tạp mà nhóm không chứng minh được là cần thiết, và phản biện Q1 sẽ đánh đúng vào đó.

\clearpage

% ================================================================= 8. TỰ PHẢN BIỆN
\section{Những chỗ đề tài còn yếu}
\label{sec:phanbien}

Mục này nhóm viết ở vai phản biện của một tạp chí Q1, cố tình tìm cách bác bỏ chính mình. Nhóm đưa vào
đề cương vì hội đồng sẽ hỏi những câu này, và tốt hơn là nhóm trả lời trước.

\begin{table}[H]\centering\small
\caption{Các điểm yếu nhóm tự xác định, mức nghiêm trọng và cách xử lý.}
\label{tab:yeu}
\begin{tabular}{@{}L{5.6cm} C{1.9cm} L{6.8cm}@{}}
\toprule
\textbf{Điểm yếu} & \textbf{Mức độ} & \textbf{Cách xử lý} \\
\midrule
Chỉ 5 sản phụ, tất cả đều đang chuyển dạ tuần 38--41. Không có điểm dữ liệu nào dưới 38 tuần, trong khi
giá trị lâm sàng chính nằm ở tuần 24--37 &
\xau{Nghiêm trọng} & Bổ sung bộ Silesia (Matonia 2020, 10 sản phụ thai kỳ) trong Phase P0, nâng $n$ lên
khoảng 22 \\
Khoảng tin cậy 95\% của 99,21 với $n = 5$ vượt quá 100\%, tức giả định phân phối chuẩn bị vi phạm &
\xau{Nghiêm trọng} & Dùng bootstrap và thang logit thay vì phân phối $t$ \\
Bảng 16 kiến trúc chạy ở dải sai, một bản ghi, một seed, ngưỡng quét trên chính bản ghi đánh giá &
\xau{Nghiêm trọng} & Chạy lại toàn bộ ở 10--60\,Hz, tách bản ghi đầy đủ, ba seed --- Phase P1 \\
Chưa cài lại baseline nào. Mọi so sánh với y văn là so với con số đã công bố &
\xau{Nghiêm trọng} & Cài lại bốn baseline cổ điển, dùng mã mở của Power-MF làm tham chiếu \\
Đóng góp C3 chưa có một dòng dữ liệu nào & Cao & Là trọng tâm Phase P4--P6 \\
Chưa chạy thí nghiệm tinh chỉnh (E7) & Trung bình & Phase P7 \\
Chưa có đường cong F1--SNR và đường cong hiệu quả mẫu & Trung bình & Phase P8 \\
Chưa tăng cường dữ liệu & Thấp & Phase P3 \\
\bottomrule
\end{tabular}
\end{table}

\subsection{Đánh giá thẳng thắn ba đóng góp}

\begin{description}[leftmargin=2.4em,style=nextline,itemsep=6pt]

\item[C2 --- đủ mạnh, là xương sống.]
Phân rã đóng góp ba tầng là kết quả phản trực giác, có kiểm định thống kê, và trực tiếp phản bác cách
làm của cả một dòng công trình. Sau khi phát biểu lại cho chính xác (\S\ref{sec:ketqua}), đây là phần
vững nhất.

\item[C3 --- tiềm năng cao nhất, nhưng chưa có dữ liệu.]
Đây là điểm khác biệt thật sự và là chỗ duy nhất nhóm còn tuyên bố ``đầu tiên'' một cách an toàn.
Kết quả ở \S\ref{sec:chonkenh} cho thấy động cơ là có thật và đã đo được. Nhưng hiện tại C3 vẫn là
thiết kế trên giấy.

\item[C1 --- yếu hơn nhóm tưởng.]
Kho \texttt{mad-lab-fau/fecg-benchmarking} đi kèm Power-MF đã tồn tại và đã làm một phần việc đối chuẩn.
Nhóm phải định vị lại C1: không phải ``đối chuẩn đầu tiên'' mà là ``đối chuẩn đầu tiên \textbf{giữa các
họ biểu diễn} ở cấu hình \textbf{đơn kênh} với quy tắc chọn kênh \textbf{mù nhãn} đăng ký trước''.

\end{description}

\subsection{Kiến trúc nhóm đã cân nhắc và loại}

Để phòng câu hỏi ``sao không dùng kiến trúc mới nhất'', nhóm đã rà soát các họ kiến trúc 2024--2026.

\begin{table}[H]\centering\small
\caption{Các họ kiến trúc mới được cân nhắc và lý do loại.}
\label{tab:loai}
\begin{tabular}{@{}L{4.2cm} C{2.0cm} L{8.0cm}@{}}
\toprule
\textbf{Họ kiến trúc} & \textbf{Quyết định} & \textbf{Lý do} \\
\midrule
Tiền huấn luyện tự giám sát kiểu wav2vec & \tot{Nên thử} & Nhóm chỉ có 5 chủ thể có nhãn, trong khi có
sẵn kho tín hiệu ổ bụng không nhãn lớn. Đây là lý do đáng tin duy nhất \\
Mamba, S4, S6 (mô hình không gian trạng thái) & Loại & Ưu thế của chúng là chuỗi rất dài. Ở đây chuỗi chỉ
1.000 mẫu và quét trường tiếp nhận đã bão hoà ở 1,5 giây \\
KAN, RetNet, xLSTM, PatchTST & Loại & Không tìm thấy căn cứ nào cho tín hiệu 4 giây ở 250\,Hz \\
Neural ODE & Loại & Chi phí lớn, không có lợi thế rõ cho bài toán định vị đỉnh \\
Transformer & Đã thử, loại & Chậm gấp 21,7 lần \texttt{cnn\_dil} để đổi lấy 0,25 điểm \\
\bottomrule
\end{tabular}
\end{table}

\begin{canhbao}[title={Ranh giới của khẳng định trên}]
Nhóm \textbf{không tìm thấy} công trình nào áp dụng Mamba cho dò QRS thai từ điện tim ổ bụng.
``Không tìm thấy'' khác với ``không tồn tại''. Nếu muốn khẳng định trong bài báo, nhóm phải tự chạy
thí nghiệm loại bỏ, chứ không được trích dẫn suy luận.
\end{canhbao}

\clearpage

% ================================================================= 9. KẾ HOẠCH
\section{Kế hoạch triển khai 24 tuần}

\subsection{Mười hai giai đoạn}

\begin{table}[H]\centering\small
\caption{Kế hoạch 12 giai đoạn trong 24 tuần, kèm hai cổng quyết định.}
\label{tab:plan}
\begin{tabular}{@{}C{1.1cm} C{1.5cm} L{6.6cm} L{4.6cm}@{}}
\toprule
\textbf{Mã} & \textbf{Tuần} & \textbf{Nội dung} & \textbf{Sản phẩm} \\
\midrule
P0 & 1--2 & Bổ sung dữ liệu: Silesia B1/B2, FECGSYNDB, NInFEA. Viết thẻ dữ liệu và quy trình chống rò rỉ
bốn tầng & Kho dữ liệu có mã băm \\
P1 & 3--4 & Chạy lại bảng kiến trúc ở 10--60\,Hz, tách bản ghi đầy đủ, ba seed & Bảng kiến trúc đã sửa \\
P2 & 5--6 & Cài lại bốn baseline cổ điển; đối chiếu với mã mở Power-MF & Baseline chạy được \\
P3 & 7--8 & Tăng cường dữ liệu; huấn luyện lại mô hình chính & Mô hình v2 \\
\midrule
\multicolumn{4}{@{}l}{\cellcolor{rowbg}\textbf{Cổng 1 (tuần 8):} nếu chưa đạt F1 $\geq$ 95\% với $n \geq 15$
chủ thể thì dừng mở rộng, tập trung viết C1 + C2} \\
\midrule
P4 & 9--11 & Cài đặt đặc trưng đồng điều bền vững cho đánh giá chất lượng & Thư viện fSQI \\
P5 & 12--13 & So fSQI topo với SampEn, bSQI, kurtosis, entropy phổ & Bảng tương quan \\
P6 & 14--16 & Cổng từ chối trả lời; đường cong rủi ro--độ phủ & Kết quả C3 \\
\midrule
\multicolumn{4}{@{}l}{\cellcolor{rowbg}\textbf{Cổng 2 (tuần 16):} nếu fSQI không đạt $|\rho| \geq 0{,}6$
thì công bố kết quả phủ định có kiểm định, vẫn đủ cho một bài báo} \\
\midrule
P7 & 17--18 & Tiền huấn luyện FECGSYNDB rồi tinh chỉnh (E7) & Kết quả chuyển miền \\
P8 & 19--20 & Đường cong F1--SNR, đường cong hiệu quả mẫu, kiểm định thống kê đầy đủ & Bộ hình cho bài báo \\
P9 & 21 & Đóng gói: API, bảng điều khiển, Docker & Sản phẩm chạy được \\
P10 & 22--23 & Viết bài hội nghị 4 trang và bản thảo tạp chí & Hai bản thảo \\
P11 & 24 & Hoàn thiện, nộp, báo cáo Euréka & Hồ sơ nộp \\
\bottomrule
\end{tabular}
\end{table}

\subsection{Phân vai ba sinh viên}

\begin{table}[H]\centering\small
\caption{Phân vai và trách nhiệm chính.}
\label{tab:vai}
\begin{tabular}{@{}L{3.2cm} L{5.4cm} L{5.6cm}@{}}
\toprule
\textbf{Vai} & \textbf{Trách nhiệm chính} & \textbf{Giai đoạn phụ trách} \\
\midrule
Chủ nhiệm đề tài & Kiến trúc mô hình, huấn luyện, thống kê, viết bài & P1, P3, P8, P10 \\
Thành viên 2 & Dữ liệu, chống rò rỉ, baseline, tái lập & P0, P2, P7 \\
Thành viên 3 & Đặc trưng topo, fSQI, cổng từ chối, đóng gói & P4, P5, P6, P9 \\
\bottomrule
\end{tabular}
\end{table}

\clearpage

% ================================================================= 10. SẢN PHẨM
\section{Sản phẩm và công bố}

\subsection{Sản phẩm bàn giao}

\begin{table}[H]\centering\small
\caption{Danh mục sản phẩm, kèm trạng thái hiện tại.}
\label{tab:sanpham}
\begin{tabular}{@{}L{5.6cm} L{6.0cm} C{2.6cm}@{}}
\toprule
\textbf{Sản phẩm} & \textbf{Mô tả} & \textbf{Trạng thái} \\
\midrule
Thư viện \texttt{fqrs\_model.py} & Cài đặt tham chiếu đầy đủ pipeline & \tot{Đã có} \\
Sáu checkpoint đã huấn luyện & Năm mô hình tách bản ghi + một mô hình sản xuất & \tot{Đã có} \\
Công cụ tải dữ liệu & Tải trực tiếp từ PhysioNet, ghi mã băm SHA-256 & \tot{Đã có} \\
Công cụ suy luận dòng lệnh & Hỗ trợ EDF, WFDB, CSV, NPY & \tot{Đã có} \\
Bộ mã đối chuẩn & Chạy đúng giao thức, kiểm chứng chéo Hungarian & \tot{Đã có} \\
Thư viện fSQI topo & Đặc trưng đồng điều bền vững & Phase P4 \\
Cổng từ chối trả lời & Kèm đường cong rủi ro--độ phủ & Phase P6 \\
API và bảng điều khiển & Đóng gói Docker & Phase P9 \\
\bottomrule
\end{tabular}
\end{table}

\subsection{Kế hoạch công bố}

\begin{enumerate}[leftmargin=1.8em,itemsep=4pt]
\item \textbf{Hội nghị Computing in Cardiology} (4 trang, hạn tháng 4): trình bày C2, phân rã đóng góp
ba tầng. Đây là kết quả đã có, viết được ngay sau Phase P1.
\item \textbf{Tạp chí Q1} --- \textit{Physiological Measurement} hoặc \textit{Biomedical Signal Processing
and Control} hoặc \textit{IEEE JBHI}: bài đầy đủ gồm C1, C2, C3.
\item \textbf{Báo cáo Euréka}: bản tiếng Việt, nhấn mạnh sản phẩm chạy được và giá trị ứng dụng.
\end{enumerate}

\begin{ghichu}[title={Vì sao đề tài có cửa Q1}]
Ba lý do, xếp theo độ vững:
\begin{enumerate}[leftmargin=1.5em,itemsep=2pt]
\item \textbf{Một kết quả phản trực giác có kiểm định thống kê} (C2): front-end tín hiệu quan trọng hơn
họ kiến trúc mạng khoảng 27 lần. Điều này phản bác cách làm của cả một dòng công trình.
\item \textbf{Một khoảng trống thật, đã kiểm tra bằng bốn cơ sở dữ liệu và ba ngôn ngữ} (G3).
\item \textbf{Một động cơ lâm sàng đã đo được} cho C3: chỉ số chất lượng dựa trên phổ đạt oracle trong
miền nhưng sụp đổ ngoài miền. Đây là bằng chứng, không phải phỏng đoán.
\end{enumerate}
Và kể cả khi C3 thất bại, giả thuyết H0 đã dự phòng: đối chuẩn cộng phân rã cộng kết quả phủ định có
kiểm định vẫn là một bài báo hoàn chỉnh.
\end{ghichu}

\clearpage

% ================================================================= 11. TÀI LIỆU
\section{Tài liệu tham khảo}

Ba mươi công trình dưới đây nhóm đã đọc \textbf{toàn văn}, không dừng ở tóm tắt. Bản PDF được lưu trong
thư mục \texttt{papers/} kèm tệp \texttt{papers\_records.json} ghi nguồn tải và mức truy cập của từng bài.

\begin{enumerate}[leftmargin=1.6em,itemsep=2pt,label={[\arabic*]}]
\item {\small Zhong W, Liao L, Guo X, Wang G. "A deep learning approach for fetal QRS complex detection." Physiological Measurement (IOP Publishing / IPEM), vol. 39, no. 4, art. 045004, 9 trang, 2018. DOI: 10.1088/1361-6579/aab297. Nhan bai 15/12/2017, sua 22/02/2018, chap nhan 27/02/2018, xuat ban 20/04/2018. Sun Yat-sen University, Guangzhou, Trung…} \hfill{\scriptsize\itshape (p01)}
\item {\small Mohebbian MR, Vedaei SS, Wahid KA, Dinh A, Marateb HR, Tavakolian K. "Fetal ECG Extraction From Maternal ECG Using Attention-Based CycleGAN." IEEE Journal of Biomedical and Health Informatics (JBHI), vol. 26, no. 2, pp. 515-526, 2022. DOI: 10.1109/JBHI.2021.3111873. University of Saskatchewan (Canada), University of Isfahan (Iran), UPC…} \hfill{\scriptsize\itshape (p03)}
\item {\small Basak P, Sakib AHMN, Chowdhury MEH, Al-Emadi N, Yalcin HC, Pedersen S, Mahmud S, Kiranyaz S, Al-Maadeed S. "A novel deep learning technique for morphology preserved fetal ECG extraction using 1D-CycleGAN." Expert Systems with Applications (Elsevier), vol. 235, art. 121196, 2024. DOI: 10.1016/j.eswa.2023.121196. University of Dhaka…} \hfill{\scriptsize\itshape (p04)}
\item {\small Arash Shokouhmand, Negar Tavassolian, "Fetal Electrocardiogram Extraction Using Dual-Path Source Separation of Single-Channel Non-Invasive Abdominal Recordings", IEEE Transactions on Biomedical Engineering, vol. 70, no. 1, pp. 283-295, January 2023. DOI: 10.1109/TBME.2022.3189617. (Nhan 09/05/2022, chinh sua 18/06/2022, chap nhan…} \hfill{\scriptsize\itshape (p05)}
\item {\small Lin Chen, Shuicai Wu, Zhuhuang Zhou, "Fetal ECG Signal Extraction from Maternal Abdominal ECG Signals Using Attention R2W-Net", Sensors, vol. 25, no. 3, art. no. 601, 2025. DOI: 10.3390/s25030601. (MDPI, open access CC BY. Nhan 22/10/2024, chinh sua 18/01/2025, chap nhan 19/01/2025, xuat ban 21/01/2025. College of Chemistry and Life…} \hfill{\scriptsize\itshape (p06)}
\item {\small Amirreza Asadi, Aboozar Ghaffari, Zahra Hatami, "U-Net-Based Deep Learning Framework for Single-Channel Fetal ECG Extraction Using Time-Delay Representation", IEEE Access, vol. 13, pp. 206367-206379+ (bai dai 14 trang), 2025. DOI: 10.1109/ACCESS.2025.3639272. (Nhan 05/11/2025, chap nhan 26/11/2025, xuat ban 03/12/2025. Biomedical…} \hfill{\scriptsize\itshape (p07)}
\item {\small Xu, Xiaojian; Zhang, Yifan; Rao, Yufei; Xu, Yinru; Gao, Yang; Tu, Huating. "A Deep Learning-Guided Ensemble Empirical Mode Decomposition Method for Single-Channel Fetal Electrocardiogram Extraction." Sensors (MDPI), 2026, tap 26, so 7, bai so 2037. DOI: 10.3390/s26072037. Nhan bai 29/01/2026, chap nhan 23/03/2026, xuat ban 25/03/2026.…} \hfill{\scriptsize\itshape (p09)}
\item {\small Esmaeili Alidash, Kimia; Danandeh Hesar, Hamed. "Deep learning-based approach for accurate detection of fetal QRS complexes in abdominal ECG signals." Scientific Reports (Nature Portfolio), 2025, tap 15, bai so 39260. DOI: 10.1038/s41598-025-22999-9. Department of Biomedical Engineering, Sahand University of Technology, Tabriz, Iran.…} \hfill{\scriptsize\itshape (p10)}
\item {\small Orvas, Iulia; Radu, Andrei; Galli, Alessandra; Neacsu, Ana; Peri, Elisabetta. "A Complex UNet Approach for Non-Invasive Fetal ECG Extraction Using Single-Channel Dry Textile Electrodes." Preprint arXiv:2506.22457v1 [eess.SP], 16/06/2025, ghi ro "preprint submitted to IEEE Journal of Biomedical and Health Informatics". Don vi: SIGMA…} \hfill{\scriptsize\itshape (p11)}
\item {\small Wahbah M, Zitouni MS, Al Sakaji R, Funamoto K, Widatalla N, Krishnan A, Kimura Y, Khandoker AH (2024). "A deep learning framework for noninvasive fetal ECG signal extraction." Frontiers in Physiology, 15:1329313. DOI: 10.3389/fphys.2024.1329313. Nhan 28/10/2023, chap nhan 22/03/2024, xuat ban 22/04/2024. Giay phep CC BY. File PDF 13…} \hfill{\scriptsize\itshape (p12)}
\item {\small Behar J (2014). "Extraction of Clinical Information from the Non-Invasive Fetal Electrocardiogram." Luan an Tien si (DPhil), Balliol College, University of Oxford. Nguoi huong dan: Prof. Gari D. Clifford (dong huong dan: Dr Julien Oster). Michaelmas 2014, 249 trang. LUU Y VE TRICH DAN: ten file trong repo ghi "arxiv" va (o ban p14) ghi…} \hfill{\scriptsize\itshape (p13)}
\item {\small Behar J (2014), Chuong 7 "Fetal QRS detection with no maternal chest reference", trang 114-137 trong luan an Tien si "Extraction of Clinical Information from the Non-Invasive Fetal Electrocardiogram", Balliol College, University of Oxford, Michaelmas 2014. Ban tap chi tuong ung (ghi trong danh muc cong bo cua chinh luan an, muc [9]…} \hfill{\scriptsize\itshape (p14)}
\item {\small Andreotti, Fernando; Behar, Joachim; Zaunseder, Sebastian; Oster, Julien; Clifford, Gari D. "An open-source framework for stress-testing non-invasive foetal ECG extraction algorithms." Physiological Measurement, tap 37, so 5, nam 2016, trang 627-648. DOI: 10.1088/0967-3334/37/5/627. (Ban PDF trong thu muc la ban luu tru CC BY-NC-ND 4.0…} \hfill{\scriptsize\itshape (p15)}
\item {\small Niknazar, Mohammad; Rivet, Bertrand; Jutten, Christian. "Fetal ECG Extraction by Extended State Kalman Filtering Based on Single-Channel Recordings." IEEE Transactions on Biomedical Engineering, tap 60, so 5, nam 2013, trang 1345-1352. DOI: 10.1109/TBME.2012.2234456. (Ban PDF trong thu muc la ban luu tru HAL, ma hal-00831928.)} \hfill{\scriptsize\itshape (p16)}
\item {\small Castillo, Encarnacion; Morales, Diego P.; Garcia, Antonio; Parrilla, Luis; Ruiz, Victor U.; Alvarez-Bermejo, Jose A. "A clustering-based method for single-channel fetal heart rate monitoring." PLoS ONE, tap 13, so 6, nam 2018, bai e0199308, 22 trang. DOI: 10.1371/journal.pone.0199308. Nhan 06/06/2017, chap nhan 05/06/2018, dang…} \hfill{\scriptsize\itshape (p17)}
\item {\small Eleonora Sulas, Monica Urru, Roberto Tumbarello, Luigi Raffo, Danilo Pani. "Systematic analysis of single- and multi-reference adaptive filters for non-invasive fetal electrocardiography". Mathematical Biosciences and Engineering (MBE), Volume 17, Issue 1, trang 286-308, 2020. DOI: 10.3934/mbe.2020016. Nhan bai 10/06/2019, chap nhan…} \hfill{\scriptsize\itshape (p19)}
\item {\small Gari D. Clifford, Ikaro Silva, Joachim Behar, George B. Moody. "Non-invasive fetal ECG analysis" (tieu de trong metadata ghi la "Noninvasive Fetal ECG analysis"). Physiological Measurement, Volume 35, Issue 8, trang 1521-1536, 2014. DOI: 10.1088/0967-3334/35/8/1521. Day la bai xa luan (editorial) mo dau so dac biet ve cuoc thi…} \hfill{\scriptsize\itshape (p20)}
\item {\small Adam Matonia, Janusz Jezewski, Tomasz Kupka, Michal Jezewski, Krzysztof Horoba, Janusz Wrobel, Robert Czabanski, Radana Kahankova. "Fetal electrocardiograms, direct and abdominal with reference heartbeat annotations". Scientific Data (Nature), volume 7, so bai 200, nam 2020. DOI: 10.1038/s41597-020-0538-z. Nhan bai 22/11/2019, chap nhan…} \hfill{\scriptsize\itshape (p21)}
\item {\small Jose A. Perea, John Harer. "Sliding Windows and Persistence: An Application of Topological Methods to Signal Analysis". File PDF trong thu muc la BAN TIEN AN arXiv:1307.6188v2 [math.AT], nop 25/11/2013, ghi ngay 22/11/2013, 34 trang; PDF KHONG ghi ten tap chi va KHONG co DOI. Ban xuat ban chinh thuc (trich duoc tu danh muc tham khao cua…} \hfill{\scriptsize\itshape (p22)}
\item {\small Henry Adams, Tegan Emerson, Michael Kirby, Rachel Neville, Chris Peterson, Patrick Shipman, Sofya Chepushtanova, Eric Hanson, Francis Motta, Lori Ziegelmeier. "Persistence Images: A Stable Vector Representation of Persistent Homology". Journal of Machine Learning Research 18 (2017) 1-35. Nop 7/2016, xuat ban 2/2017. Bien tap vien:…} \hfill{\scriptsize\itshape (p23)}
\item {\small Mathieu Carriere, Frederic Chazal, Yuichi Ike, Theo Lacombe, Martin Royer, Yuhei Umeda. "PersLay: A Neural Network Layer for Persistence Diagrams and New Graph Topological Signatures". Proceedings of the 23rd International Conference on Artificial Intelligence and Statistics (AISTATS) 2020, Palermo, Italy. PMLR Volume 108. Ban quyen…} \hfill{\scriptsize\itshape (p24)}
\item {\small Anirudh Som, Hongjun Choi, Karthikeyan Natesan Ramamurthy, Matthew P. Buman, Pavan Turaga. "PI-Net: A Deep Learning Approach to Extract Topological Persistence Images". Metadata cua file PDF ghi subject = "IEEE Conference on Computer Vision and Pattern Recognition Workshops" (CVPRW), nam 2020. Ma nguon:…} \hfill{\scriptsize\itshape (p25)}
\item {\small Renata Turkes (University of Antwerp), Guido Montufar (UCLA), Nina Otter (Queen Mary University of London). "On the Effectiveness of Persistent Homology". arXiv:2206.10551v3 [math.AT], ban v3 ngay 16 thang 1 nam 2023 (ban PDF day du kem phu luc, 39 trang). Ban rut gon 17 trang trong cung thu muc (p26\_Turkes2022-EffectivenessPH.pdf) co…} \hfill{\scriptsize\itshape (p26)}
\item {\small Quoc Hoan Tran, Yoshihiko Hasegawa (Department of Information and Communication Engineering, Graduate School of Information Science and Technology, The University of Tokyo). "Topological time-series analysis with delay-variant embedding". arXiv:1803.00208v4 [physics.data-an], ban v4 ghi ngay 10-11 thang 1 nam 2019. Ma nguon:…} \hfill{\scriptsize\itshape (p27)}
\item {\small Eugene Tan, Shannon Algar, Débora Corrêa, Michael Small, Thomas Stemler, David Walker. "Selecting embedding delays: An overview of embedding techniques and a new method using persistent homology." arXiv:2302.03447v1 [math.DS], nộp 7-2-2023, trang đầu ghi ngày 8-2-2023. Complex Systems Group, Department of Mathematics and Statistics, The…} \hfill{\scriptsize\itshape (p28)}
\item {\small Yonglian Ren, Feifei Liu, Shengxiang Xia, Shuhua Shi, Lei Chen, Ziyu Wang. "Dynamic ECG signal quality evaluation based on persistent homology and GoogLeNet method." Frontiers in Neuroscience, tập 17, số bài 1153386, 2023. DOI: 10.3389/fnins.2023.1153386. Nhận 29-1-2023, chấp nhận 20-2-2023, xuất bản 8-3-2023. Đơn vị: School of Science,…} \hfill{\scriptsize\itshape (p29)}
\item {\small Meryll Dindin, Yuhei Umeda, Frédéric Chazal. "Topological Data Analysis for Arrhythmia Detection through Modular Neural Networks." arXiv:1906.05795v1 [cs.LG], 13-6-2019. Đơn vị: École Centrale Paris (Dindin), Fujitsu Nhật Bản (Umeda), INRIA (Chazal). File PDF này KHÔNG ghi DOI và KHÔNG ghi tên hội nghị. LƯU Ý: trong danh mục tài liệu…} \hfill{\scriptsize\itshape (p30)}
\item {\small Yu-Min Chung, Chuan-Shen Hu, Yu-Lun Lo, Hau-Tieng Wu. "A Persistent Homology Approach to Heart Rate Variability Analysis with an Application to Sleep-Wake Classification". Ban trong file la tien an pham arXiv:1908.06856v2 [eess.SP], 02/05/2020, 38 trang. FILE PDF NAY KHONG IN TEN TAP CHI VA KHONG IN DOI. Ban xuat ban chinh thuc duoc…} \hfill{\scriptsize\itshape (p31)}
\item {\small Andy Dominguez-Monterroza, Alfonso Mateos Caballero, Antonio Jimenez-Martin (Department of Artificial Intelligence, Universidad Politecnica de Madrid, Tay Ban Nha). "Age-dependent patterns of cardiac complexity unveiled by topological data analysis of pediatric heart rate variability". PLoS One. 2025;20(12):e0337620. doi:…} \hfill{\scriptsize\itshape (p32)}
\item {\small Alperen Karan, Atabey Kaygun (Department of Mathematical Engineering, Istanbul Technical University, 34467 Istanbul, Tho Nhi Ky). "Time Series Classification via Topological Data Analysis". arXiv:2102.01956v2 [stat.ML], 11/06/2021, 31 trang. Email: karana@itu.edu.tr, kaygun@itu.edu.tr. FILE PDF NAY KHONG IN TEN TAP CHI VA KHONG IN DOI -…} \hfill{\scriptsize\itshape (p33)}
\end{enumerate}

\subsection*{Tài liệu bổ sung phát hiện trong quá trình rà soát}

\begin{enumerate}[leftmargin=1.6em,itemsep=2pt,label={[\arabic*]},start=31]
\item Jaeger KM, Nissen M, Rahm S, Titzmann A, Fasching PA, Beilner J, Eskofier BM, Leutheuser H.
Power-MF: robust fetal QRS detection from non-invasive fetal electrocardiogram recordings.
\textit{Physiological Measurement} 45(5):055009, 2024. DOI 10.1088/1361-6579/ad4952.
Mã nguồn: \texttt{github.com/mad-lab-fau/fecg-benchmarking}.
\end{enumerate}

\clearpage

% ================================================================= 12. PHỤ LỤC
\section*{Phụ lục A --- Cách chạy lại toàn bộ kết quả}
\addcontentsline{toc}{section}{Phụ lục A --- Cách chạy lại toàn bộ kết quả}

\begin{verbatim}
pip install numpy scipy scikit-learn torch mne wfdb

# 1. Tải dữ liệu trực tiếp từ PhysioNet, ghi mã băm SHA-256
cd model
python download_data.py --root data --only adfecgdb

# 2. Huấn luyện lại toàn bộ (khoảng 30 phút CPU), sinh 6 checkpoint
python train_final.py --epochs 6 --seed 0

# 3. Suy luận trên một bản ghi bằng checkpoint chưa từng thấy bản ghi đó
python predict.py --input data/adfecgdb/r01.edf --lead 1 --annot qrs \
                  --checkpoint checkpoints/fetalqrs_tcn_fold_r01.pt

# 4. Chạy đối chuẩn và quy tắc chọn kênh mù nhãn
cd ../benchmark_dpss
python run_dpss_protocol.py --mode loro
python full_measure.py
python blind_lead.py
\end{verbatim}

\section*{Phụ lục B --- Cấu trúc thư mục bàn giao}
\addcontentsline{toc}{section}{Phụ lục B --- Cấu trúc thư mục bàn giao}

\begin{table}[H]\centering\small
\begin{tabular}{@{}L{5.0cm} L{9.4cm}@{}}
\toprule
\textbf{Thư mục} & \textbf{Nội dung} \\
\midrule
\texttt{model/} & Thư viện tham chiếu, mã huấn luyện, mã suy luận, 6 checkpoint, dữ liệu ADFECGDB \\
\texttt{pilot\_evidence/} & Mã và nhật ký của toàn bộ thí nghiệm tiền khả thi \\
\texttt{benchmark\_dpss/} & Mã đối chuẩn, kiểm chứng Hungarian, quy tắc chọn kênh mù nhãn \\
\texttt{papers/} & 30 bản PDF toàn văn và ma trận bằng chứng \\
\texttt{survey/} & Dữ liệu khảo sát 30 công trình, tệp số liệu đã kiểm chứng \\
\texttt{de\_cuong\_latex/} & Mã nguồn LaTeX của đề cương này, hình vector, mã sinh bảng \\
\bottomrule
\end{tabular}
\end{table}

\vfill
\begin{ghichu}[title={Cam kết về tính liêm chính}]
\small
Mọi con số trong đề cương này đều do nhóm tự đo và chạy lại được. Nơi nào nhóm chưa chắc chắn,
nhóm ghi rõ. Nơi nào bản đề cương trước đó của nhóm phát biểu sai, nhóm nêu rõ chỗ sai và con số đúng
(\S\ref{sec:ketqua}, \S\ref{sec:chonkenh}). Nhóm không tuyên bố ``vượt trội'' ở những chỗ giao thức
không cho phép so sánh trực tiếp.
\end{ghichu}
